\section{Dataset Evaluation}
\label{sec:evaluation}

\subsection{Hypothesis}
\label{sec:hypothesis}

The new Elixir typechecker is designed to be more error-prone than Dialyzer.
% Dialyzer is built around success typing: it warns only when it can prove that a discrepancy is certain, accepting many function definitions whose behaviours are merely suspicious.
The set-theoretic typechecker performs a structural analysis of each function body against its declared and inferred types and warns on definitions that violate them.
On this setting, every function definition Dialyzer warns should also be warned by the typechecker, and the typechecker should additionally find potential type errors Dialyzer cannot detect.
Formally:
\begin{gather*}
  \text{Dialyzer fails} \;\Rightarrow\; \text{Typechecker fails} \\
  \text{Dialyzer passes} \;\Rightarrow\; \text{Typechecker fails $\vee$ passes}
\end{gather*}
We expect the implications to hold together with the strict superset expectation that the typechecker fails on substantially more definitions than Dialyzer.

To judge whether the hypothesis holds we examine the \emph{intersection} of the two tools' warnings, that is the fraction of both warned by Dialyzer and the typechecker over only Dialyzer-warned.
The closer the fraction to one confirms the hypothesis: the new system preserves Dialyzer's discrimination and exceeds it.
Where the fraction yields less than one, we then reason about the residual, i.e., the entries Dialyzer warns on but the typechecker does not.

The evaluation addresses three questions:
\begin{enumerate}
    \item To what extent is the typechecker stricter than Dialyzer, and how much of Dialyzer's discrimination does it subsume?
    \item Where the subsumption breaks down, why does Dialyzer catch an error that the typechecker does not, and is that a fundamental disagreement or a limitation of the present state of either the typechecker or the translation?
    \item What does the new typechecker detect that Dialyzer does not, and what does this imply for downstream use?
\end{enumerate}

\subsection{Overall Results and Verdict}
\label{sec:overall-results}

The dataset construction described in Section~\ref{sec:dataset-construction} produced \num{23073} translated function specifications.
Dialyzer ran to completion on all of them; the new typechecker analysed \num{22621} and could not analyse the remaining 452, whose host projects fail to compile under the custom compiler (a dependency that the development-branch compiler rejects aborts the whole project; see Section~\ref{sec:dataset-construction}).
The cross-comparison below is therefore over the \num{22621} entries both tools resolved, with the 452 reported separately.
Table~\ref{tab:overall-results} summarises the outcomes:

\begin{table}[h]
    \centering
    \begin{tabular}{lrr}
        \toprule
        \textbf{Outcome} & \textbf{Count} & \textbf{\% of total} \\
        \midrule
        Both tools passed                   & \num{18688} & 81.0\% \\
        \midrule
        Typechecker-only warned             & \num{3636} & 15.8\% \\
        Dialyzer-only warned                &     98   &  0.4\% \\
        Both tools warned                   &    199   &  0.9\% \\
        \midrule
        Typechecker could not analyse       &    452   &  2.0\% \\
        \midrule
        Dialyzer passed (total)             & \num{22685} & 98.3\% \\
        Typechecker passed (total)          & \num{18786} & 81.4\% \\
        \bottomrule
    \end{tabular}
    \caption{Tool outcomes across \num{23073} translated function specifications. The typechecker rejects \num{3835} entries (16.6\%) against Dialyzer's 297 (1.3\%); a further 452 it could not analyse.}
    \label{tab:overall-results}
\end{table}

Two observations stand out.
First, the typechecker is decisively the stricter of the two analyses: it rejects \num{3835} entries (16.6\% of the corpus) where Dialyzer rejects only 297 (1.3\%), a thirteen-fold difference.
The set-theoretic typechecker's pass rate (81.4\%) sits below Dialyzer's (98.3\%).
Second, where the two tools fail, they overlap substantially but not completely.

\paragraph{Subsumption rate on failures.}
Dialyzer warned on 297 entries in total; the typechecker also warned on 199 of them, giving a subsumption rate on failures of $199/297 \approx 67.0\%$.
The remaining 98 Dialyzer-warned entries (the Dialyzer-only set) are the cases where the first implication appears to break down.
However, as Subsection~\ref{sec:dialyzer-only} establishes, none of these 98 represents a case in which the typechecker disagrees with Dialyzer; each is attributable to an analysis the typechecker does not perform, to a feature (\texttt{@opaque}-type enforcement) it does not yet model, or to precision lost during translation.

\paragraph{Verdict.}
The data support the hypothesis's subsumption clause but qualify its strict-superset reading.
The new typechecker does warn on far more definitions than Dialyzer, just as predicted; but, as Subsection~\ref{sec:typecheck-only} shows, most of that excess is translation residuals --- struct patterns the translator could not fully resolve from the AST and positions widened to \texttt{dynamic()} --- rather than superior error detection.
The first implication, that if Dialyzer warns the typechecker warns too, holds directly for 67.0\% of Dialyzer's warnings, and none of the remaining cases is a real disagreement.
Of the 98 entries Dialyzer catches but the typechecker does not, each is attributable to a whole-program check the typechecker does not perform, a finding it makes locally but that translation masked with \texttt{dynamic()}, or an \texttt{@opaque}-abstraction violation the type system does not yet model, none a genuine type-theoretic disagreement (Subsection~\ref{sec:dialyzer-only} gives the breakdown).
In short, the two analyses are complementary rather than nested: where both resolve the same types the typechecker reproduces Dialyzer's findings, but each also flags a class the other cannot, namely Dialyzer the whole-program properties and intra-project calls the typechecker leaves as \texttt{dynamic()}, and the typechecker a core of purely-static body-level return leaks that Dialyzer's optimism accepts (Subsection~\ref{sec:typecheck-only}).

\subsection{Approximation via \texttt{dynamic()}}
\label{sec:approximation-via-dynamic}

Out of the \num{23073} entries, \textbf{\num{13910} (60.3\%)} contain at least one \texttt{dynamic()} sub-expression in the translated annotation.
This widening is the single most consequential approximation in the translation, and it is the mechanism behind a large share of the residual disagreement analysed below.

When, for example, the definition of a remote type \texttt{Module.type\_name()} is not available at translation time (because the module lives in a dependency outside the translation path), the translator conservatively replaces the reference with \texttt{dynamic()}.
Because \texttt{dynamic()} is compatible with every type, the resulting annotation is no stronger than the original specification with respect to its unresolved components: the typechecker, checking a body against an argument typed \texttt{dynamic()}, cannot derive a contradiction that Dialyzer derives from the original concrete contract.

A subtlety of the new system sharpens this picture: \texttt{dynamic()} is always hoisted to the root of its enclosing container~\cite{elixir-set-theoretic-types}.
When the translation produces a nested gradual type such as \texttt{\{:ok, dynamic()\}}, Elixir rewrites it to \texttt{dynamic(\{:ok, term()\})}: the whole tuple becomes gradual, not merely its second element.
The widening is thus coarser than its textual position suggests: a single unresolved remote type anywhere inside a structure makes the entire structure gradual, so the masking analysed below operates at the granularity of the enclosing tuple, map, or list rather than the individual leaf.
The same property carries an important converse: an annotation containing \emph{no} \texttt{dynamic()} at all is checked as a fully static type: ``if the user provides their own types, and those types are not \texttt{dynamic()}, then Elixir's type system behaves as a statically typed one''~\cite{elixir-set-theoretic-types}.
A dynamic-free annotation is therefore a genuinely static, fully checked specification, which is exactly why the precise training pool (Subsection~\ref{sec:downstream-implications}) and the safe-and-precise criterion require \texttt{dynamic()} to be \emph{entirely absent} rather than merely rare.
The direct consequence is the one to keep in mind when reading a \emph{pass}: where the annotation carries \texttt{dynamic()} the typechecker accepts trivially, so a pass certifies only the concrete positions and leaves open whether the static type would pass or warn. For a wholly \texttt{dynamic()} annotation, such as the opacity cases of Subsection~\ref{sec:dialyzer-only}, a pass is no static verdict at all.

The prevalence of \texttt{dynamic()} does not invalidate the translation; it quantifies how much of the corpus carries cross-module type references that cannot be refined without expanding the translation scope to the full dependency tree.
Its effect runs in both directions and is made precise in the next two subsections: it suppresses some Dialyzer-detectable errors (Subsection~\ref{sec:dialyzer-only}), yet \num{2747} of the \num{3835} typechecker-rejected entries (72\%) still carry a \texttt{dynamic()} and are rejected regardless, showing that the widening is far from a blanket suppression of the typechecker.

\subsection{Dialyzer-Only Warnings: Why Dialyzer Catches What the Typechecker Cannot \emph{Yet}}
\label{sec:dialyzer-only}

The 98 Dialyzer-only entries are the crux of the hypothesis: they are precisely the cases where Dialyzer rejects a definition the typechecker accepts.
However, establishing that none of them is a true type-theoretic counterexample is what allows the first implication to stand.
Table~\ref{tab:dialyzer-cats} breaks them down by warning category; the category names and their meanings follow Dialyzer's own classification~\cite{dialyzer-manual}, as surfaced through the \texttt{dialyxir} interface used during construction~\cite{dialyxir}.
A counting caveat first: Dialyzer routinely reports several warnings for a single function definition (these 98 entries carry 134 warnings between them), whereas the new typechecker generally surfaces only one per function, a difference that matters once both tools fire on the same definition (Subsection~\ref{sec:both-warned}).
The column \textbf{N} below therefore counts the \emph{entries} that carry each category, not raw warnings.
A representative entry for each category, namely the function and the Dialyzer message it draws, none of which the typechecker reproduces, is given in the appendix ``\nameref{app:dialyzer-examples}''.

\begin{table}[h]
\centering
\footnotesize
\begin{tabularx}{\textwidth}{@{}l r X l@{}}
\toprule
\textbf{Category} & \textbf{N} & \textbf{What Dialyzer reports} & \textbf{In typechecker?} \\
\midrule
\texttt{no\_return}            & 29 & A function never returns normally; it always raises, exits, or loops. & no (whole-program) \\
\texttt{pattern\_match}        & 18 & A pattern can never match the type of the term it is matched against. & partial (local) \\
\texttt{call\_without\_opaque} & 12 & An \texttt{@opaque}-typed argument is supplied by its raw internal structure. & no (opacity) \\
\texttt{unused\_fun}           & 12 & A function is defined but never called anywhere. & no (whole-program) \\
\texttt{call\_to\_missing}     & 11 & A call targets a function missing or not exported by its module. & no (cross-module) \\
\texttt{unknown\_function}     &  9 & A call targets a function that does not exist. & no (cross-module) \\
\texttt{pattern\_match\_cov}   &  5 & A clause is already fully covered by earlier clauses and is unreachable. & partial (local) \\
\texttt{guard\_fail}           &  4 & A guard test can never succeed for the inferred types. & partial (local) \\
\texttt{call}                  &  3 & A call's arguments cannot match the callee's expected input types, so it can never succeed. & partial (local) \\
\texttt{callback\_type\_mismatch} & 1 & A \texttt{@behaviour} callback implementation's spec conflicts with the behaviour's declared callback type. & no (cross-module) \\
\texttt{opaque\_match}         &  1 & A pattern inspects the internal structure of an \texttt{@opaque} term. & no (opacity) \\
\bottomrule
\end{tabularx}
\caption{Dialyzer warning categories among the 98 Dialyzer-only entries. \textbf{N} is the number of entries carrying each category; an entry with warnings in several categories appears in more than one row, so the column sums to more than 98. The final column indicates whether the typechecker performs an analysis that can produce the same finding. \emph{partial (local)} marks a category the typechecker checks inside a function body and therefore reproduces when both tools resolve the same types (Subsection~\ref{sec:both-warned}), but misses here because translation widened the relevant type to \texttt{dynamic()}; \emph{no} marks a finding outside its analysis altogether, requiring \emph{whole-program} reasoning (call-graph reachability or dead code), \emph{cross-module} resolution, or an \emph{opacity} model it does not have.}
\label{tab:dialyzer-cats}
\end{table}

Classifying the 98 entries by \emph{why} the typechecker does not reproduce the warning yields a clean three-way partition:
\begin{itemize}
  \item \textbf{59 entries} fall solely in categories the set-theoretic typechecker does not analyse at all: the whole-program reachability and dead code behind \texttt{no\_return} and \texttt{unused\_fun}, and the cross-module resolution behind \texttt{unknown\_function}, \texttt{call\_to\_missing}, and \texttt{callback\_type\_mismatch}.
  \item \textbf{27 entries} carry a finding the typechecker performs \emph{locally} and reproduces elsewhere, namely an argument mismatch (\texttt{call}), a dead pattern or clause (\texttt{pattern\_match}, \texttt{pattern\_match\_cov}), or an impossible guard (\texttt{guard\_fail}), but missed here; the causes, translation widening the scrutinised type to \texttt{dynamic()} versus a genuine residual gap, are examined below.
  \item the remaining \textbf{12 entries} carry \emph{solely} an opacity warning (\texttt{call\_without\_opaque}, \texttt{opaque\_match}) for \texttt{@opaque} type that the typechecker does not model and the translation cannot represent.
\end{itemize}
None of the three is a case where the typechecker, given the same information and feature set, type-theoretically disagrees with Dialyzer. We treat each in turn.

\paragraph{Whole-program analyses outside its scope (59 entries).}
By design the new typechecker performs \emph{module type inference}: it infers and checks each function using the current module, the standard library, and dependencies, while assuming that calls to other modules \emph{within the same project} have type \texttt{dynamic()}; its inference is best-effort, warning only on definite failures (cases where \emph{every} combination of a type fails), so a passed function is \emph{not disproven} rather than verified sound~\cite{elixir-set-theoretic-types}.
Five Dialyzer categories fall outside this body-local view entirely.
\texttt{no\_return} and \texttt{unused\_fun} are products of Dialyzer's \emph{whole-program} success-typing fixed point~\cite{lindahl2006practical}: detecting that a function has no reachable return path, or that it is never called, requires reasoning across the program's control flow and call sites, which a current function-at-a-time typechecker does not do.
\texttt{unknown\_function}, \texttt{call\_to\_missing}, and \texttt{callback\_type\_mismatch} require \emph{cross-module} resolution: because the typechecker assumes any call to another module in the same project has type \texttt{dynamic()}, it never looks the callee up, and so never notices that the target function is missing or unexported (\texttt{unknown\_function}, \texttt{call\_to\_missing}) or that an implementation violates the callback type its \texttt{@behaviour} declares (\texttt{callback\_type\_mismatch}).
On these 59 entries the typechecker found no body-level fault either, so it passed; the warnings are therefore not type disagreements but findings of an analysis it does not \emph{yet} perform, and nothing in its theory precludes adding them.
Until then Dialyzer remains complementary here, and the 59 entries are evidence of an analytical-scope difference, not of a definition the typechecker wrongly accepts.

A further consideration discounts the \texttt{unknown\_function} and \texttt{call\_to\_missing} warnings in particular as trivial.
Dialyzer emits one whenever a called function's module is absent from its PLT (the persistent lookup table of analysed modules), so the warning reports \emph{dependency-closure completeness}, not type correctness: the residual cases here call into optional OTP applications (\texttt{:eldap}), development and test tooling (\texttt{IEx}, \texttt{Mix}, \texttt{ExUnit}), optional libraries outside the build closure (\texttt{Jason}, \texttt{PhoenixHTMLHelpers}), or compiler-internal and operator-expansion helpers (\texttt{:elixir\_quote}, \texttt{Enum.\_\_in\_\_/2}) that no project declares as an analysed dependency.
This sensitivity to the PLT is also why the count is environment-dependent (a handful of these warnings in earlier constructions disappeared once the relevant dependencies were built into the PLT), and it confirms the category is an artefact of the build closure rather than a property of the code.
The typechecker, which never resolves cross-project calls in the first place, has no analogue to produce, so these entries cannot be a disagreement.

\paragraph{Locally checkable but unreproduced (27 entries).}
The remaining non-opacity categories are all within the typechecker's body-local remit: an argument mismatch at a call site (\texttt{call}), a pattern or clause that can never match (\texttt{pattern\_match}, \texttt{pattern\_match\_cov}), and a guard that can never succeed (\texttt{guard\_fail}) are each a body-against-declaration inconsistency it is built to catch.
That it does catch them is visible in the both-warned set (Subsection~\ref{sec:both-warned}), where it independently emits the same finding on 36 of 78 \texttt{call} entries, 30 of 55 \texttt{pattern\_match}, and a few of the others.
That it did not reproduce these 27 is a matter of approximation and analysis scope, not of an inability to express the check itself.
For only 6 of them could the cause be direct \texttt{dynamic()} masking: the annotation still carries a \texttt{dynamic()} that can cover the position Dialyzer faults, so where Dialyzer reasons against the original concrete type and finds the call impossible or the pattern dead, the widened annotation presents the typechecker with a value that vacuously satisfies the call and makes every pattern reachable (Subsection~\ref{sec:approximation-via-dynamic}); with the type translated concretely the typechecker would be expected to reproduce the finding.
The other 21 carry \emph{no} \texttt{dynamic()} at all, and their annotations are faithful renderings (a source \texttt{map()} becomes \texttt{open\_map()}, an \texttt{any()} becomes \texttt{term()}, with nothing widened), so they are a genuine residual gap: real dead-clause, unreachable-pattern, and impossible-guard findings the typechecker simply does not make.
Each of the 21 hinges on whole-program information the module-local analysis lacks. The unreachability follows from the return type of an intra-project call (a same-project callee returns an error tuple where the code expects a map, as in \texttt{as\_to\_model\_data/1}'s \texttt{extract\_actors/1} and the Mobilizon and Akkoma federation fetchers, which module-local inference assumes is \texttt{dynamic()}); from caller-side success typing that narrows an argument below its declared type (a four-atom tag only ever \texttt{:hour} or \texttt{:week} in \texttt{ttl\_to\_seconds/1}); or from a compile-time module attribute the typechecker does not constant-fold (\texttt{if @env == :test}, statically dead in \texttt{task\_children/0} and \texttt{maybe\_inserted\_at/0}).
All need whole-program success typing or compile-time constant-folding that the module-local, best-effort typechecker does not perform: the Dialyzer categories are ones the typechecker can express, but these particular instances require information beyond its scope. They are the same whole-program reach that yields the 59 above, surfacing here under a locally-checkable category rather than as a translation artefact.

\paragraph{Opacity warnings the typechecker does not model (12 entries).}
The remaining entries carry \texttt{opaque\_match} or \texttt{call\_without\_opaque}: Dialyzer warns that code constructs or pattern-matches an \texttt{@opaque} type through its internal structure rather than its public interface, an abstraction-boundary violation, here on library types such as \texttt{Ecto.Multi} and \texttt{Mint.HTTP.t()}, and even a project-local opaque (\texttt{Concentrate.TripDescriptor}).
This is not a type error in the ordinary sense but an opacity-enforcement finding, and the new type system does not distinguish an opaque type from its underlying structure: it does not model the \texttt{@opaque} (including \texttt{@typep}) type at all.
The translation cannot supply what the target system lacks: an opaque type is bluntly translated to its structure, or, when its defining dependency is unavailable, to \texttt{dynamic()}, hence the boundary of type security levels is erased on both sides.
Empirically the widening applies to eleven of the twelve: they carry a \texttt{dynamic()} in the translated annotation and none names the opaque type at all.
Crucially, and unlike the locally-checkable group, this is \emph{not} a precision loss that resolving the dependency would repair: even with the opaque type fully available, the typechecker would still not flag the violation, because it has no notion of opacity to violate.
Tellingly, the one entry whose annotation contains no \texttt{dynamic()} at all (\texttt{ExCubicIngestion.ProcessIngestion.ingest/1}, on \texttt{Ecto.Multi}) is missed all the same---confirming that opacity, not approximation, is the operative cause.

\paragraph{Project distribution.}
The Dialyzer-only warnings are concentrated: Mobilizon contributes 50 of the 98 and Akkoma 20, with the remaining 28 spread across twelve further projects (MBTA seven; Analytics five; Absinthe and Cachex three each; cldr and Livebook two each; and six projects with one apiece).
The concentration tracks \texttt{dynamic()} density: the projects with the most unresolved remote types are those whose Dialyzer findings the translation most often masks with the gradual type.

\subsection{Typechecker-Only Warnings}
\label{sec:typecheck-only}

\num{3636} entries triggered a warning from the typechecker that Dialyzer did not report.
At face value this is the strict-superset clause of the hypothesis; on inspection, however, most of it is not superior error detection but translation residuals and the conservatism of module-local inference.
The findings divide by kind into unreachable clauses or patterns (\num{1677}), argument mismatches at a call site (914), return-type mismatches (691), map-key accesses (217), and a residue of other checks (137).

\paragraph{Unreachable-clause findings (\num{1677}).}
The largest group is dominated by a translation artefact rather than genuine dead code: \num{1308} of the \num{1677} reject a struct pattern the translated type cannot satisfy, and every one is a residue of translating structs from the AST without a symbol table (Subsection~\ref{sec:struct-representation}).
In 752 the translated struct type lists only some of the struct's fields, so the closed map excludes valid structs and a \texttt{\%Mod\{\}} pattern is rejected; in 556 the translator resolved the struct through a short or aliased name it could not expand to the full module path (\texttt{\%ViewPort\{\}} for \texttt{Scenic.ViewPort}), so the asserted \texttt{\_\_struct\_\_} atom differs from the pattern's.
Neither is a type-theoretic disagreement: both are the cost of AST-based translation without a persistent lookup table, and both would be removed by expanding the translation scope to the defining modules.
The remaining 369 are non-struct, and they include the genuine dead code the typechecker surfaces and Dialyzer accepts: a clause matching an atom outside the scrutinee's type (\texttt{:dev\_blue} where the value is only \texttt{:dev\_green} or \texttt{:prod}), a defensive fall-through made redundant by an exhaustive earlier match, a right-hand \texttt{||} branch that can never run because the left side is always truthy, or a constantly-true-or-false guard.
For a codebase moving off Dialyzer onto the new checker, this reachability and exhaustiveness feedback is precisely the kind of static signal the transition is meant to add.

\paragraph{Argument mismatches (914).}
These are the conservatism of module-local inference: the faulted argument is inferred \texttt{term()} or \texttt{dynamic()}, from an intra-project call or a broad declared type such as \texttt{map()} or \texttt{list()}, so the typechecker cannot establish the callee's expected input and warns on correct code; others pass a struct whose translated type dropped the field the callee reads, turn on protocol dispatch (for instance \texttt{Ecto.Queryable}) the checker does not model, or are spec-looseness, where a bare \texttt{map()} parameter becomes an open map that guarantees no field, so a field access in the body cannot be shown to succeed even though the runtime map carries the key.

\paragraph{Return-type findings (691).}
These classify by the typechecker's inferred body type against the declared type into three groups.
In \num{390} the annotation itself carries a \texttt{dynamic()}: after root-hoisting it is gradual, so the disagreement is on a non-static type and neither a pass nor a warning here is a static verdict (Subsection~\ref{sec:approximation-via-dynamic}); this group absorbs the translation artefacts, a concrete body warned against a \texttt{dynamic()} assertion that root-hoisting stripped of its optionality, or over-broadened inference that invents a union arm the body never returns.
In a further 125 the inferred return type is \texttt{dynamic()} or the top type \texttt{term()} while the asserted type is concrete: the body returns the result of an intra-project or un-spec'd remote call that inference assumes is \texttt{dynamic()}, so the typechecker cannot establish the concrete return type and warns, whereas Dialyzer's whole-program success typing resolves the callee and confirms the contract (\texttt{def config\_path, do: Path.join(data\_dir(), \dots)} is inferred \texttt{dynamic()} against a declared \texttt{binary()} because \texttt{data\_dir/0} is intra-project).
The remaining 176 are purely static, no \texttt{dynamic()} on either side, so a disagreement is a true static verdict: the body provably returns a value outside the declared type that Dialyzer's optimism accepts, as when \texttt{Scenic.PubSub.publish/2} returns \texttt{\{:error, :not\_registered\}} under the specification \texttt{:ok} or \texttt{Absinthe.Type.Custom.parse\_date/1} returns \texttt{:error} under \texttt{\{:ok, nil\}}.
These purely-static leaks are the legitimate content of the strict-superset clause; representative cases are given in the appendix ``\nameref{app:typecheck-examples}''.

\paragraph{Interpretation.}
The headline ``\num{3636} warnings Dialyzer misses'' overstates the typechecker's advantage.
The unreachable-clause group is dominated by struct-representation residuals of AST-based translation (\num{1308} of \num{1677}), not error detection, and the return-type group by \texttt{dynamic()}-masked and cannot-verify cases that carry no static verdict (515 of 691).
What the typechecker alone genuinely catches is the 176 purely-static return leaks Dialyzer's optimism accepts and the 369 non-struct unreachable-clause findings its reachability checking surfaces.
The two analyses are thus better described as \emph{complementary} than nested for now: Dialyzer flags the whole-program properties and intra-project calls the typechecker leaves as \texttt{dynamic()}, and the typechecker both the return leaks Dialyzer's optimism accepts and the dead code its reachability checking surfaces.
As the migration target shifts from Dialyzer to the new checker, these two classes, the return leaks and the reachability findings, are what the codebase stands to gain.

\subsection{Entries Flagged by Both Tools}
\label{sec:both-warned}

199 entries triggered warnings from both tools, the core of the subsumption rate.
Their Dialyzer categories are led by \texttt{call} (78), \texttt{no\_return} (70), \texttt{pattern\_match} (55), and \texttt{call\_without\_opaque} (24), a mix of the locally-checkable and whole-program categories of Subsection~\ref{sec:dialyzer-only}, but here arising in functions whose bodies \emph{also} violate their declared return or argument types.
These are the unambiguous errors: a definition wrong enough that Dialyzer's success typing and the typechecker's body check independently reject it.
The typechecker reports a single warning per function (158 of the 199 carry exactly one), and which one it is depends on the category: for the locally-checkable categories it reproduces Dialyzer's finding directly (incompatible arguments on 36 of 78 \texttt{call} entries, a dead pattern on 30 of 55 \texttt{pattern\_match}), whereas for the whole-program categories it cannot, and rejects instead on a coexisting body-level fault; across the 35 entries carrying only whole-program or cross-module Dialyzer categories, that fault is an unreachable clause in 21, a map-key access in 6, an argument mismatch in 5, a return-type mismatch in 2, and another check in 1.
Either way the entry is co-rejected: the typechecker need not match Dialyzer's diagnosis, only fail the function on some ground.
Their existence confirms that the two analyses agree wherever a definition's fault is visible to both, and that the disagreement studied above is confined to the analytical-scope and approximation boundaries, not to the verdict on plainly faulty code.
One entry per Dialyzer category from this set, paired with both tools' messages, is given in the appendix ``\nameref{app:both-examples}''.

\subsection{Implications for Downstream Use}
\label{sec:downstream-implications}

The two-tool cross-validation is not only a quality check on the translation; it directly shapes how the resulting dataset is consumed in the LLM training pipeline of the next chapter, whose goal is a model that predicts annotations the typechecker accepts, that is, predicts \emph{safe}.
Three observations transfer directly.

\paragraph{Keeping only typechecker-verified labels.}
The model is trained to produce annotations the typechecker accepts, so every training label must itself be an annotation the typechecker admits.
We therefore keep only the \emph{both-pass} entries, those accepted by \emph{both} Dialyzer and the new typechecker.
This is a substantive filter, not a cosmetic one: it removes \num{4385} entries (19.0\% of the corpus), the great majority typechecker rejections, plus the 452 the typechecker could not analyse and the 297 Dialyzer warnings.
About a fifth of the translated annotations are thus rejected by, or unverifiable under, at least one tool (genuinely inconsistent for some, but for most only unverifiable by the typechecker's conservative inference; see Subsection~\ref{sec:typecheck-only}), and excluding them leaves only labels the typechecker positively accepts.

\paragraph{Two datasets: with and without gradual types.}
Within the both-pass pool one choice remains: whether to also drop entries whose annotation contains \texttt{dynamic()}.
Subsection~\ref{sec:approximation-via-dynamic} showed this \texttt{dynamic()} is mostly an artefact of unresolved remote types, not a deliberate gradual-typing decision, and the choice has costs on both sides: dropping it yields purely static, precise labels but discards more than half the pool, while keeping it preserves coverage at the risk of teaching the model to emit \texttt{dynamic()} to avoid type errors.
Rather than decide in advance, we build \emph{two} datasets and let the next chapter settle the matter empirically:
\begin{itemize}
  \item the full \textbf{both-pass} pool (\textbf{\num{18688}} entries), which \emph{includes} gradual (\texttt{dynamic()}) types; and
  \item its \textbf{dynamic-free} subset (\textbf{\num{7813}} entries) of \emph{purely static} types only.
\end{itemize}
The two are nested (the second is the precise core of the first); a length filter and a subproject-stratified split (Subsection~\ref{sec:llm-data-prep}) give the final per-track counts.
Training one model on each lets us test directly whether including gradual types helps or hurts safe, precise prediction.

\paragraph{Typechecker as the success criterion.}
Because the model is optimised to satisfy the typechecker, evaluation uses the typechecker as ground truth: a predicted annotation counts as correct when the typechecker accepts it on injection into the real project.
Dialyzer plays no part in scoring predictions: its role was confined to building the dataset (the agreement study of this chapter and the both-pass filter). The prediction metric is the typechecker's acceptance alone, which is appropriate because the task is precisely to produce annotations the new type system admits.

\subsection{Summary}

\begin{itemize}
    \item \textbf{Coverage}: all \num{23073} specs from 383 repositories were translated; Dialyzer analysed every one, and the new typechecker all but 452 whose host projects fail to compile under the development-branch compiler, leaving \num{22621} with a verdict from both tools.
    \item \textbf{Relative strictness}: the typechecker rejects \num{3835} entries (16.6\%) against Dialyzer's 297 (1.3\%); pass rates are 81.4\% versus 98.3\%. The typechecker is decisively the stricter tool, as hypothesised.
    \item \textbf{Subsumption}: the typechecker independently warns on 67.0\% of Dialyzer's failing entries (199 of 297), reproducing Dialyzer's finding directly for the locally-checkable categories and co-rejecting the rest on a separate body-level fault. The 98 it does not co-warn are \emph{not} type-theoretic counterexamples, but split between whole-program analyses it does not perform, locally-checkable findings masked by \texttt{dynamic()}, and \texttt{@opaque} type violations it does not model (Subsection~\ref{sec:dialyzer-only}).
    \item \textbf{Approximation cost}: 60.3\% of entries contain \texttt{dynamic()}; it suppresses some Dialyzer-visible errors but leaves 72\% of typechecker rejections intact, so it is not a blanket suppression.
    \item \textbf{Typechecker-only (15.8\%)}: its dominant kind is unreachable-clause findings, of which \num{1308} of \num{1677} reject struct patterns the translator could not fully resolve from the AST (incomplete field sets or unresolved module paths) rather than genuine errors; the return-type group is likewise dominated by \texttt{dynamic()}-masked and cannot-verify cases carrying no static verdict. What it alone genuinely catches is 176 purely-static return leaks Dialyzer's optimism accepts and 369 non-struct dead-clause findings its reachability checking surfaces (Subsection~\ref{sec:typecheck-only}).
    \item \textbf{Both tools (0.9\%)}: definitions faulty enough that both analyses independently reject them, confirming agreement wherever a fault is visible to both.
    \item \textbf{Training pools}: the both-pass filter yields \num{18688} typechecker-verified entries; its dynamic-free subset (\num{7813}) and the full set define the two supervision tracks compared empirically in the next chapter.
\end{itemize}
The hypothesis is supported in its subsumption clause but qualified in its strict-superset clause.
The set-theoretic typechecker subsumes Dialyzer wherever both reason over the same resolvable types. Its residual blind spots on Dialyzer's failures are attributable to whole-program analyses it has not \emph{yet} implemented, to features such as \texttt{@opaque}-type enforcement it does not yet model, and to translation-time approximation, not to any case where it wrongly accepts a definition Dialyzer soundly rejects.
It also warns far more often than Dialyzer, but that excess is mostly translation residuals --- struct patterns it could not fully resolve from the AST and positions widened to \texttt{dynamic()} --- rather than superior error detection, with a genuine core of 176 purely-static return leaks and 369 non-struct dead-clause findings that it alone catches.
The two analyses are therefore best read as complementary: each detects a class the other cannot.
